\section{Recomendaciones de la entrevista y preguntas frecuentes}
\subsection{Repaso de las preguntas del público}
Graham respondió varias preguntas frecuentes: ¿puede el Agent encargarse de tareas que no son de código? ¿cómo se mide la tasa de éxito del Agent? ¿cuándo hace falta que un humano tome el control? Sus respuestas reflejan la convicción central de la «colaboración humano-máquina»: el Agent se encarga de la parte repetible y orquestable, y el humano conserva la autoridad de revisión y control de riesgos.

\begin{itemize}
    \item «¿Puede el Agent hacer análisis de requisitos?» → Puede ayudar a organizar el texto, pero el PM todavía tiene que indagar los puntos ciegos.
    \item «¿Cómo se toma el control cuando algo falla?» → El diff que genera el Agent lo revisa un human in the loop, y si no se aprueba se revierte.
    \item «¿Cómo se conserva la traceability?» → OpenHands registra todas las acciones, y el registro incluye la marca de tiempo, el comando y el resultado.
\end{itemize}

\subsection{Lista de recomendaciones}
\begin{table}[H]
\centering
\begin{tabular}{p{0.33\textwidth} p{0.63\textwidth}}
\toprule
Tema & Recomendación \\
\midrule
Alineación de requisitos & Dividir los requisitos en tareas pequeñas y especificar la salida esperada durante la etapa de Plan. \\
\midrule
Habilitación de registros & Hacer que el Agent escriba «intención + resultado» en el registro, para facilitar el análisis humano. \\
\midrule
Control de seguridad & Configurar las operaciones de base de datos y los comandos de despliegue como pasos que requieren approve humano. \\
\bottomrule
\end{tabular}
\caption{Panorama rápido de las recomendaciones frecuentes.}
\end{table}

\subsection{Resumen de la sección}
\begin{itemize}
    \item El Q\&A de la entrevista vuelve a subrayar que el Agent no sustituye al ingeniero humano, sino que es su compañero de apoyo.
    \item Una lista de recomendaciones clara facilita que el equipo lleve a la práctica el contenido de la entrevista.
\end{itemize}

